% Chapter 4: System Design (Expanded for 80+ Page Project)
\chapter{SYSTEM DESIGN}

\section{System Architecture Overview}
The Nebula Multi-Cloud platform is designed using a multi-layered, asynchronous architecture. This design ensures that the platform remains responsive, even when performing long-running cloud provisioning tasks like spinning up a virtual machine or a large database.

\subsection{Cloud-Centric Presentation Layer}
The user interface is a "State-Driven React 19 Application." It communicates with the backend via a secure, RESTful API. The key requirement for the presentation layer was modern usability—it allows users to "Cockpit" their entire multi-cloud estate from a single screen.

\subsection{Asynchronous Backend (FastAPI + Celery)}
Traditional Web servers (like Flask or Django in sync mode) would time out during a Terraform "Apply" operation, which can take anywhere from 2 to 10 minutes. Nebula solves this by using **FastAPI** as the API controller and **Celery** as the background worker farm. When a user clicks "Provision," the API returns a "Task ID" immediately, and the Celery worker handles the actual infrastructure execution in the background.

\section{Entity-Relationship (ER) Design and SQL Schemas}
The database (PostgreSQL) is the "Brain" of the platform, storing cloud inventory, credentials, and user sessions.

\subsection{Core Tables Description}
\begin{enumerate}
    \item \textbf{User Table}: Stores hashed passwords and authentication metadata.
    \item \textbf{Cloud Credential Table}: Stores the encrypted (AES-128) JSON blobs for AWS, Azure, and GCP keys.
    \item \textbf{Project Table}: Groups resources into logical business units.
    \item \textbf{Resource Table}: Stores the state, provider IDs, and configuration for every provisioned cloud object.
\end{enumerate}

\section{Data Flow Diagrams (DFD)}
System data flow was modeled at three levels of abstraction:
\begin{itemize}
    \item \textbf{Level 0 (Context)}: User provides "Cloud API Credentials" and receives "Infrastructure Status" and "Log Remediation."
    \item \textbf{Level 1 (Orchestration)}: Internal flow between the API, the Security Vault, the Redis task queue, and the Cloud Providers.
    \item \textbf{Level 2 (Deep Dive)}: Detailed flow of the "Provisioning Engine" during a Terraform Init-Plan-Apply cycle.
\end{itemize}

\newpage
% Chapter 5: Implementation (Expanded for 80+ Page Project)
\chapter{IMPLEMENTATION}

\section{Technology Stack Deep-Dive}
We chose a modern, cutting-edge stack to handle the complexities of multi-cloud management.

\subsection{FastAPI: High-Performance Backend}
FastAPI was chosen for its high-throughput (ASGI) capabilities and integrated support for Pydantic. Pydantic ensures that data received from the React frontend is strictly typed and validated before it reaches the "Provisioning Engine."

\subsection{Terraform: The Universal Cloud Translator}
Instead of writing millions of lines of code to call AWS and Azure SDKs, Nebula uses **Terraform** as its execution backend. Nebula dynamically generates a "Workspace" folder for each task, populates it with HCL (HashiCorp Configuration Language) files, and executes the provider-specific API calls.

\section{Detailed Creation Steps for Cloud Resources}
This section details exactly how a user creates resources on the Nebula platform. This workflow is consistent across storage, compute, and networking.

\subsection{Steps for Provisioning Virtual Machines (AWS/Azure/GCP)}
1. \textbf{Task Selection}: User selects "Create VM" and chooses the desired provider.
2. \textbf{Configuration Form}: 
    \begin{itemize}
        \item \textit{Name}: Give the instance a unique identifier.
        \item \textit{Image}: Select the OS (e.g., Ubuntu 22.04 LTS).
        \item \textit{Instance Size}: Choose the vCPU and RAM (e.g., t3.medium).
    \end{itemize}
3. \textbf{Credential Verification}: The backend fetches the encrypted AWS/Azure key from the vault.
4. \textbf{Provisioning Trigger}: Nebula starts the background Celery task.
5. \textbf{Log Persistence}: As Terraform runs, Nebula streams the STDOUT/STDERR logs directly back to the user's screen in real-time.

\subsection{Steps for Creating Storage Buckets}
1. \textbf{S3/Blob/GCS Choice}: Select the object storage provider.
2. \textbf{Bucket Naming}: Input a globally unique name (checked for errors locally).
3. \textbf{Access Level}: Choose from "Private," "Public," or "Public-Read."
4. \textbf{Lifecycle Rules}: (Optional) Define when data should be moved to "Cold Storage."
5. \textbf{Apply}: Nebula initiates the bucket creation API through the state-manager.

\section{Core Algorithm: The AI Copilot}
The AI Copilot does not just "guess"—it uses a fine-tuned LLM model that was trained on thousands of Terraform "Apply" failures. 
\begin{enumerate}
    \item The system captures the "stderr" output from a failed task.
    \item It strips out any sensitive identifiers like IP addresses or secret keys.
    \item The scrubbed log is sent to the "Nebula AI Assistant."
    \item The assistant interprets errors like "InsufficientPermissions" or "InvalidInstanceID" and provides a "1-Click Fix."
\end{enumerate}

\newpage
